\documentclass[12pt, a4paper]{article}
\usepackage[utf8]{inputenc}
\usepackage[T1]{fontenc}
\usepackage{hyperref}
\usepackage{geometry}
\usepackage{listings}
\usepackage{xcolor}
\geometry{margin=2.5cm}

\title{\textbf{Dokumentasi Perubahan DapurPintar: \\ Implementasi Fitur Dark Mode dan Halaman Admin}}
\author{
  Raditya Rafif Rizqullah (NPM 51423172) \\
  [Nama Penulis 2] (NPM [NPM Penulis 2]) \\
  [Nama Penulis 3] (NPM [NPM Penulis 3]) \\
  [Nama Penulis 4] (NPM [NPM Penulis 4])
}
\date{5 Mei 2026}

\begin{document}

\maketitle
\thispagestyle{empty}
\newpage

\section{Pendahuluan}
DapurPintar adalah sebuah platform berbasis web yang dirancang untuk membantu pengguna dalam mengelola resep masakan dan perencanaan menu harian menggunakan kecerdasan buatan. Website ini dapat diakses secara publik melalui \url{https://dapur-pintar-psnq5xlhs-radityajulius-projects.vercel.app/}, sementara repositori kode sumber terbaru tersimpan di \texttt{C:/Users/Raditya/Documents/GitHub/DapurPintar/dapur-pintar}.

Tugas ini mendokumentasikan pembaruan terbaru yang dilakukan pada repositori pada tanggal 5 Mei 2026, yaitu penambahan fitur \textbf{Dark Mode} dan \textbf{Halaman Admin}, sebagai bagian dari tugas kuliah Informatika Universitas Gunadarma.

\section{Fitur Dark Mode}

\subsection{Ringkasan}
Fitur Dark Mode ditambahkan ke seluruh halaman aplikasi DapurPintar menggunakan library \texttt{next-themes} dan Tailwind CSS v4. Fitur ini memungkinkan pengguna mengalihkan tampilan antara mode terang (light) dan mode gelap (dark) secara manual, atau mengikuti pengaturan sistem operasi.

\subsection{Implementasi Teknis}

\subsubsection{Instalasi Package}
\begin{lstlisting}[language=bash]
npm install next-themes
\end{lstlisting}

\subsubsection{Konfigurasi CSS (globals.css)}
Penambahan variabel CSS untuk dark mode dengan mendefinisikan class \texttt{.dark} yang berisi variabel warna alternatif:
\begin{itemize}
  \item \texttt{--background: \#0a0a0a} (latar belakang gelap)
  \item \texttt{--foreground: \#ededed} (teks terang)
  \item \texttt{--card: \#171717} (kartu gelap)
  \item \texttt{--primary: \#00cc88} (warna utama versi gelap)
\end{itemize}

Menggunakan \texttt{@custom-variant dark (\&:where(.dark, .dark *))} untuk mendukung Tailwind v4.

\subsubsection{Theme Provider (ThemeProvider.tsx)}
File baru \texttt{src/app/ThemeProvider.tsx} dibuat untuk membungkus aplikasi dengan \texttt{NextThemesProvider}:
\begin{lstlisting}[language=tsx]
"use client";
import { ThemeProvider as NextThemesProvider } from "next-themes";

export default function ThemeProvider({
  children,
}: {
  children: React.ReactNode;
}) {
  return (
    <NextThemesProvider attribute="class" defaultTheme="system" enableSystem>
      {children}
    </NextThemesProvider>
  );
}
\end{lstlisting}

\subsubsection{Integrasi pada Layout dan Komponen}
\begin{itemize}
  \item \textbf{layout.tsx}: Membungkus \texttt{children} dengan \texttt{<ThemeProvider>} dan menambahkan \texttt{suppressHydrationWarning}
  \item \textbf{Navbar.tsx}: Menambahkan tombol toggle dengan ikon \texttt{Sun} dan \texttt{Moon} dari \texttt{lucide-react}
  \item \textbf{page.tsx (Landing Page)}: Mengupdate semua elemen dengan \texttt{dark:} variants
  \item \textbf{login/page.tsx dan register/page.tsx}: Menyesuaikan form dengan mode gelap
\end{itemize}

\section{Fitur Halaman Admin}

\subsection{Ringkasan}
Halaman Admin ditambahkan untuk memungkinkan pengelola website memantau statistik penggunaan aplikasi, termasuk jumlah user, recipe yang dihasilkan, dan data lainnya secara real-time.

\subsection{Perubahan Database}
Penambahan field pada model \texttt{User} di \texttt{prisma/schema.prisma}:
\begin{lstlisting}
isAdmin Boolean @default(false)
\end{lstlisting}

Migrasi database dijalankan dengan perintah:
\begin{lstlisting}[language=bash]
npx prisma migrate dev --name add_is_admin
\end{lstlisting}

\subsection{API Routes Baru}

\subsubsection{POST /api/admin/login}
\begin{itemize}
  \item Menerima email dan password
  \item Memverifikasi user dan mengecek flag \texttt{isAdmin: true}
  \item Mengembalikan token admin dengan payload JWT \texttt{\{ userId, isAdmin: true \}}
\end{itemize}

\subsubsection{GET /api/admin/stats}
Mengembalikan statistik lengkap:
\begin{itemize}
  \item \texttt{totalUsers}, \texttt{totalRecipes}, \texttt{totalSavedRecipes}
  \item \texttt{usersThisWeek}, \texttt{recipesThisWeek}
  \item \texttt{recipesByMood}, \texttt{recipesByMealType}, \texttt{recipesByLanguage}
  \item \texttt{userGrowth} (6 bulan terakhir)
  \item \texttt{recentUsers}, \texttt{recentRecipes}
\end{itemize}

\subsection{Halaman Admin}

\subsubsection{Admin Layout (src/app/admin/layout.tsx)}
\begin{itemize}
  \item Proteksi route: mengecek \texttt{adminToken} di localStorage
  \item Redirect otomatis ke \texttt{/admin/login} jika belum terautentikasi
  \item Navbar dengan tombol logout
\end{itemize}

\subsubsection{Admin Login Page (src/app/admin/login/page.tsx)}
\begin{itemize}
  \item Form login dengan email dan password
  \item POST ke \texttt{/api/admin/login}
  \item Menyimpan token di localStorage sebagai \texttt{adminToken}
  \item Redirect ke \texttt{/admin/dashboard} jika sukses
  \item Mendukung dark mode
\end{itemize}

\subsubsection{Admin Dashboard Page (src/app/admin/dashboard/page.tsx)}
Menampilkan:
\begin{itemize}
  \item Cards statistik: Total Users, Total Recipes, Saved Recipes, New This Week
  \item Charts: Recipes by Mood, Recipes by Meal Type, User Growth (6 bulan)
  \item Tabel: Recent Users (5 terbaru), Recent Recipes (5 terbaru)
  \item Mendukung dark mode
\end{itemize}

\subsection{Komponen Pendukung}
\textbf{StatsCard.tsx}: Komponen card statistik reusable dengan dukungan dark mode (icon, title, value, description).

\subsection{Keamanan}
\begin{itemize}
  \item Admin token menggunakan JWT dengan payload \texttt{\{ userId, isAdmin: true \}}
  \item Semua API \texttt{/api/admin/*} memverifikasi token dan flag \texttt{isAdmin}
  \item Token disimpan di localStorage (rekomendasi: gunakan HttpOnly cookies untuk produksi)
\end{itemize}

\section{Struktur File yang Diubah/Dibuat}

\begin{tabular}{|l|l|l|}
\hline
\textbf{No} & \textbf{File} & \textbf{Status} \\
\hline
1 & \texttt{prisma/schema.prisma} & Diubah \\
2 & \texttt{lib/auth.ts} & Diubah (generateAdminToken) \\
3 & \texttt{src/app/api/auth/me/route.ts} & Diubah (include isAdmin) \\
4 & \texttt{src/app/api/admin/login/route.ts} & \textbf{BARU} \\
5 & \texttt{src/app/api/admin/stats/route.ts} & \textbf{BARU} \\
6 & \texttt{src/components/admin/StatsCard.tsx} & \textbf{BARU} \\
7 & \texttt{src/app/admin/layout.tsx} & \textbf{BARU} \\
8 & \texttt{src/app/admin/login/page.tsx} & \textbf{BARU} \\
9 & \texttt{src/app/admin/dashboard/page.tsx} & \textbf{BARU} \\
10 & \texttt{src/app/ThemeProvider.tsx} & \textbf{BARU} \\
11 & \texttt{src/app/globals.css} & Diubah (dark mode vars) \\
12 & \texttt{src/app/layout.tsx} & Diubah (ThemeProvider) \\
13 & \texttt{src/components/Navbar.tsx} & Diubah (toggle button) \\
14 & \texttt{src/app/page.tsx} & Diubah (dark variants) \\
\hline
\end{tabular}

\section{Repositori dan Website}
Berikut adalah tautan terkait proyek DapurPintar:
\begin{itemize}
  \item Website: \url{https://dapur-pintar-psnq5xlhs-radityajulius-projects.vercel.app/}
  \item Repositori GitHub: \texttt{C:/Users/Raditya/Documents/GitHub/DapurPintar/dapur-pintar}
  \item Build Status: \texttt{npm run build} - Compiled successfully
\end{itemize}

\section{Kesimpulan}
Dokumentasi ini merangkum perubahan yang dilakukan pada proyek DapurPintar pada tanggal 5 Mei 2026. Pembaruan meliputi:
\begin{enumerate}
  \item \textbf{Dark Mode}: Implementasi menggunakan next-themes dan Tailwind v4 dengan toggle button di navbar
  \item \textbf{Halaman Admin}: Sistem autentikasi admin, API statistik, dan dashboard lengkap dengan visualisasi data
\end{enumerate}

Kedua fitur ini telah diuji dengan \texttt{npm run build} dan berhasil dikompilasi tanpa error. Fitur dark mode meningkatkan kenyamanan visual pengguna, sementara halaman admin memberikan kontrol penuh bagi pengelola website untuk memantau aktivitas pengguna.

\end{document}